\documentclass[a4paper,10pt]{report}

\def\packagepath{../../preambule}   % path du package princial
\usepackage{\packagepath/preambule}    % utilisation du fichier de configuration

\def\level{Première }              % Classe
\def\course{NSI}              % Matière
\def\eval{DS04 - CORRECTION}

\def\visibleornot{visible}    % visible or invisible
\def\documentpath{./Sources_Latex}
\def\date{Jeudi 06 Octobre 2025}
\renewcommand{\arraystretch}{1}  % Ecart dans les tableau

\def\ptsexoA{2}
\def\ptsexoB{1}
\def\ptsexoC{1}
\def\ptsexoD{1}
\def\ptsexoE{1}
\def\ptsexoF{1}
\def\ptsexoG{2}
\def\ptsexoH{1}
\def\ptsexoI{2}
\def\ptsexoJ{2}
\def\ptsexoK{1}

\def\ptstotal{\ptsexoA+\ptsexoB}

\usepackage{hyperref}

\usetikzlibrary{shapes, arrows, positioning}
\usetikzlibrary{shapes.geometric, arrows, positioning, decorations.pathreplacing}

\tikzstyle{etat} = [draw, rounded corners=15pt, minimum width=2.5cm, minimum height=1.2cm, text centered, font=\bfseries]
\tikzstyle{nouveau} = [etat, fill=blue!40]
\tikzstyle{pret} = [etat, fill=green!60]
\tikzstyle{elu} = [etat, fill=yellow!60]
\tikzstyle{bloque} = [etat, fill=red!60]
\tikzstyle{termine} = [etat, fill=white]
\tikzstyle{fleche} = [thick, ->, >=stealth]

\usepackage{float} % Permet d'utiliser [H]
\usepackage[T1]{fontenc}
\begin{document}

\renewcommand{\labelitemi}{\textbullet}

\pagestyle{DS_FP}          % Feuille de style fancy
\NomPrenomNote{}


\exods{\ptsexoA}

Sans aucune justification, donner le nombre d'états possibles codables avec:

\begin{enumerate}[1)]
    \item 1 bit \dotfill{} \textcolor{red}{\textbf{2}}
    \item 3 bits \dotfill{} \textcolor{red}{\textbf{8}}
    \item 8 bits \dotfill{} \textcolor{red}{\textbf{256}}
    \item n bits \dotfill{} \textcolor{red}{\textbf{$2^n$}}
\end{enumerate}

\medskip


\begin{minipage}{0.49\linewidth}
 \exods{\ptsexoB}

   Convertir le nombre binaire suivant en décimal. Vous justifierez vos réponses.

\begin{answer}{1}{5} 
    \begin{center}
        \verb|1010 0011|
    \end{center}   
\end{answer}

\textcolor{red}{
\textbf{Réponse :}
\begin{align*}
1010\,0011_2 &= 1 \times 2^7 + 0 \times 2^6 + 1 \times 2^5 + 0 \times 2^4 \\
&\quad + 0 \times 2^3 + 0 \times 2^2 + 1 \times 2^1 + 1 \times 2^0 \\
&= 128 + 32 + 2 + 1 \\
&= \textbf{163}_{10}
\end{align*}
}
\end{minipage}\hfill
\begin{minipage}{0.49\linewidth}
     \exods{\ptsexoC}

    Convertir le nombre décimal suivant en binaire. Vous justifierez vos réponses.
\begin{answer}{1}{5} 
    \begin{center}
        \verb|159|
    \end{center}   
\end{answer}

\textcolor{red}{
\textbf{Réponse :} Divisions successives par 2
\begin{align*}
159 \div 2 &= 79 \text{ reste } \textbf{1} \\
79 \div 2 &= 39 \text{ reste } \textbf{1} \\
39 \div 2 &= 19 \text{ reste } \textbf{1} \\
19 \div 2 &= 9 \text{ reste } \textbf{1} \\
9 \div 2 &= 4 \text{ reste } \textbf{1} \\
4 \div 2 &= 2 \text{ reste } \textbf{0} \\
2 \div 2 &= 1 \text{ reste } \textbf{0} \\
1 \div 2 &= 0 \text{ reste } \textbf{1}
\end{align*}
En lisant de bas en haut : \textbf{10011111}$_2$
}
\end{minipage}

\medskip


\begin{minipage}{0.49\linewidth}
 \exods{\ptsexoD}

   Convertir le nombre hexadécimal suivant en décimal. Vous justifierez vos réponses.

\begin{answer}{1}{5} 
    \begin{center}
        \verb|AF|
    \end{center}   
\end{answer}

\textcolor{red}{
\textbf{Réponse :}
\begin{align*}
AF_{16} &= A \times 16^1 + F \times 16^0 \\
&= 10 \times 16 + 15 \times 1 \\
&= 160 + 15 \\
&= \textbf{175}_{10}
\end{align*}
}
\end{minipage}\hfill
\begin{minipage}{0.49\linewidth}
     \exods{\ptsexoE}

Convertir le nombre binaire suivant en hexadécimal. Vous justifierez vos réponses.
\begin{answer}{1}{5} 
    \begin{center}
        \verb|1110 0110|
    \end{center}   
\end{answer}

\textcolor{red}{
\textbf{Réponse :} Regroupement par 4 bits
\begin{align*}
1110\,0110_2 &\rightarrow \text{groupes : } 1110 \text{ et } 0110 \\
1110_2 &= 8 + 4 + 2 = 14 = E_{16} \\
0110_2 &= 4 + 2 = 6_{16} \\
&\Rightarrow \textbf{E6}_{16}
\end{align*}
}
\end{minipage}

\medskip

\exods{\ptsexoF}

Effectuer l'addition binaire suivante:



    \begin{answer}{1}{4} \begin{center}
        \verb|10110111+1101110|
\end{center}  \end{answer}

\textcolor{red}{
\textbf{Réponse :}
\begin{verbatim}
    1 1 1  1         (retenues)
    1 0 1 1 0 1 1 1
  +   1 1 0 1 1 1 0
  -----------------
  1 0 0 1 0 0 1 0 1
\end{verbatim}
\textbf{Résultat : 100100101}$_2$

Vérification : $183 + 110 = 293_{10} = 100100101_2$ \checkmark
}

\pagebreak

\exods{\ptsexoG}

Expliquer simplement comment déterminer si un nombre binaire est pair ou impair. Justifier.

    \begin{answer}{1}{3}    \end{answer}

\textcolor{red}{
\textbf{Réponse :}

\textbf{Méthode :} Il suffit de regarder le \textbf{bit de poids faible} (le dernier bit, à droite).

\begin{itemize}
    \item Si le dernier bit est \textbf{0} $\rightarrow$ le nombre est \textbf{pair}
    \item Si le dernier bit est \textbf{1} $\rightarrow$ le nombre est \textbf{impair}
\end{itemize}

\textbf{Justification :} En binaire, chaque position représente une puissance de 2. Toutes les puissances de 2 (sauf $2^0 = 1$) sont paires. Donc :
\begin{itemize}
    \item Si le bit de poids faible est 0, la somme ne contient que des multiples de 2 $\rightarrow$ nombre pair
    \item Si le bit de poids faible est 1, on ajoute 1 à une somme paire $\rightarrow$ nombre impair
\end{itemize}

\textbf{Exemples :}
\begin{itemize}
    \item $1010_2 = 10_{10}$ (pair) $\rightarrow$ dernier bit = 0
    \item $111_2 = 7_{10}$ (impair) $\rightarrow$ dernier bit = 1
\end{itemize}
}



\medskip

\exods{\ptsexoH}

Ecrire une fonction \verb|est_pair(n)| qui prend en paramètre un nombre écrit en binaire (au format chaine de caractères) et renvoie \verb|True| si le nombre est pair et \verb|False| sinon.

    \begin{answer}{1}{7} 
\begin{verbatim}
def est_pair(n):
    ''' Renvoie True si n est pair et False sinon
    n chaîne de caractères représentant un nombre binaire
    sortie: bool '''









assert est_pair('1010') == True    # 10 en décimal
assert est_pair('111') == False    # 7 en décimal
    \end{verbatim}

    \end{answer}

\textcolor{red}{
\textbf{Correction :}
\begin{verbatim}
def est_pair(n):
    ''' Renvoie True si n est pair et False sinon
    n: chaîne de caractères représentant un nombre binaire
    sortie: bool '''
    return n[-1] == '0'
\end{verbatim}
\textbf{Explication :} On vérifie simplement si le dernier caractère de la chaîne est '0'.
}

\medskip

\exods{\ptsexoI}

Ecrire une fonction \verb|nb_bits(n)| qui prend en paramètre un entier naturel \verb|n| et renvoie le nombre de bits nécessaires pour l'écrire en binaire. Par exemple, \verb|nb_bits(5)| renvoie \verb|3| car $5 = 101_2$.

On pourra par exemple compter le nombre de divisions entières par 2 nécessaires pour obtenir 0.
    \begin{answer}{1}{9.2} 
\begin{verbatim}
def nb_bits(n):
    ''' Renvoie le nombre de bits nécessaires pour écrire n en binaire 
    n entier naturel (int)
    sortie: int '''














assert nb_bits(5) == 3    # 5 = 101_2
assert nb_bits(10) == 4   # 10 = 1010_2 
assert nb_bits(16) == 5   # 16 = 10000_2
    \end{verbatim}

    \end{answer}

\textcolor{red}{
\textbf{Correction :}
\begin{verbatim}
def nb_bits(n):
    ''' Renvoie le nombre de bits nécessaires pour écrire n en binaire
    n: entier naturel (int)
    sortie: int '''
    if n == 0:
        return 1
    
    compteur = 0
    while n > 0:
        n = n // 2
        compteur += 1
    
    return compteur
\end{verbatim}
\textbf{Explication :} On divise successivement par 2 jusqu'à obtenir 0, en comptant le nombre d'itérations. Cas particulier : 0 nécessite 1 bit.
}

    \pagebreak

    \exods{\ptsexoJ}

Ecrire une fonction \verb|nb_uns(val_bin)| qui prend en paramètre une chaine de caractères \verb|val_bin| représentant un nombre binaire et qui renvoie le nombre de bits à 1 dans cette chaine.

    \begin{answer}{1}{9}
\begin{verbatim}
def nb_uns(val_bin):
    ''' Renvoie le nombre de bits à 1 dans val_bin
    val_bin chaine de caractères représentant un nombre binaire
    













assert nb_uns('101010') == 3
assert nb_uns('1111') == 4
assert nb_uns('0000') == 0
    \end{verbatim}

    \end{answer}

\textcolor{red}{
\textbf{Correction :}
\begin{verbatim}
def nb_uns(val_bin):
    ''' Renvoie le nombre de bits à 1 dans val_bin
    val_bin: chaine de caractères représentant un nombre binaire
    sortie: int '''
    return val_bin.count('1')
\end{verbatim}
\textbf{Explication :} La méthode \verb|.count()| compte le nombre d'occurrences du caractère '1' dans la chaîne.

\textbf{Alternative avec boucle :}
\begin{verbatim}
def nb_uns(val_bin):
    compteur = 0
    for bit in val_bin:
        if bit == '1':
            compteur += 1
    return compteur
\end{verbatim}
}

    \medskip

     \exods{\ptsexoK}

Ecrire une fonction \verb|multiplication(val_bin, puissance)| qui va multiplier un nombre binaire \verb|val_bin| (sous forme de chaine de caractères) par $2^{puissance}$ où \verb|puissance| est un entier naturel.
    \begin{answer}{1}{10}
\begin{verbatim}
def multiplication(val_bin, puissance):
    ''' Multiplie le nombre binaire val_bin par 2^puissance
    val_bin chaine de caractères représentant un nombre binaire
    puissance entier naturel (int)
    sortie: str '''














assert multiplication('101', 3) == '101000'    # 5 * 2^3 = 40
assert multiplication('111', 2) == '11100'     # 7 * 2^2 = 28
assert multiplication('1001', 0) == '1001'     # 9 * 2^0 = 9
    \end{verbatim}

    \end{answer}

\textcolor{red}{
\textbf{Correction :}
\begin{verbatim}
def multiplication(val_bin, puissance):
    ''' Multiplie le nombre binaire val_bin par 2^puissance
    val_bin: chaine de caractères représentant un nombre binaire
    puissance: entier naturel (int)
    sortie: str '''
    return val_bin + '0' * puissance
\end{verbatim}
\textbf{Explication :} Multiplier par $2^n$ en binaire revient à décaler le nombre vers la gauche de n positions, c'est-à-dire ajouter n zéros à droite.

\textbf{Exemples :}
\begin{itemize}
    \item $101_2 \times 2^3 = 101000_2$ (on ajoute 3 zéros)
    \item $111_2 \times 2^2 = 11100_2$ (on ajoute 2 zéros)
\end{itemize}
}

    \pagebreak
    \thispagestyle{DS_LP}

\exods{0}

    On dispose de deux chaines binaires \verb|A| et \verb|B| de même longueur (8 bits).

    On veut produire une chaine binaire représentant leur somme.

    Ecrire une fonction \verb|addition_binaire(A, B)| qui prend en paramètres les deux chaines binaires \verb|A| et \verb|B| et renvoie la chaine binaire représentant leur somme. On ne pourra pas convertir les valeurs en décimal avant de faire l'addition.

\begin{answer}{1}{11} 
\begin{verbatim}
def addition_binaire(A, B):
    ''' Renvoie la chaine binaire représentant la somme de A et B
    A et B chaines de caractères représentant des nombres binaires de même longueur (8 bits)
    sortie: str '''















assert addition_binaire('00011010', '00000101') == '00011111'    # 26 + 5 = 31
assert addition_binaire('11111111', '00000001') == '100000000'    # 255 + 1 = 256
assert addition_binaire('10101010', '01010101') == '11111111'    # 170 + 85 = 255
    \end{verbatim}

    \end{answer}

\textcolor{red}{
\textbf{Correction :}
\begin{verbatim}
def addition_binaire(A, B):
    ''' Renvoie la chaine binaire représentant la somme de A et B
    A et B: chaines de caractères représentant des nombres binaires 
            de même longueur (8 bits)
    sortie: str '''
    resultat = ""
    retenue = 0
    
    # On parcourt de droite à gauche
    for i in range(len(A) - 1, -1, -1):
        bit_a = int(A[i])
        bit_b = int(B[i])
        
        somme = bit_a + bit_b + retenue
        
        resultat = str(somme % 2) + resultat
        retenue = somme // 2
    
    # S'il reste une retenue à la fin
    if retenue == 1:
        resultat = '1' + resultat
    
    return resultat
\end{verbatim}

\textbf{Explication de l'algorithme :}
\begin{enumerate}
    \item On parcourt les bits de droite à gauche (indices décroissants)
    \item Pour chaque position, on additionne les deux bits + la retenue précédente
    \item Le bit résultat est \verb|somme % 2| (reste de la division par 2)
    \item La nouvelle retenue est \verb|somme // 2| (quotient de la division par 2)
    \item Si une retenue finale existe, on l'ajoute à gauche
\end{enumerate}

\textbf{Table d'addition binaire avec retenue :}
\begin{verbatim}
0 + 0 + 0 = 0 (retenue 0)    0 + 1 + 1 = 0 (retenue 1)
0 + 0 + 1 = 1 (retenue 0)    1 + 0 + 0 = 1 (retenue 0)
0 + 1 + 0 = 1 (retenue 0)    1 + 0 + 1 = 0 (retenue 1)
1 + 1 + 0 = 0 (retenue 1)    1 + 1 + 1 = 1 (retenue 1)
\end{verbatim}
}

\end{document}